\documentclass{article}
\usepackage[margin=2cm]{geometry}
\usepackage{amsmath}
\usepackage{graphicx}
\usepackage{booktabs}
\usepackage[utf8]{inputenc}
\usepackage{ragged2e}
\setlength{\emergencystretch}{3em}
\begin{document}
\sloppy

\noindent\hfill\includegraphics[height=1.8em]{assets/figure_header_logo.png}\par\noindent\hrule\vspace{4pt}

and considering the empirical distribution function Fn(), the halfspace
depthwillbe


\begin{equation}
d _ { H S } ( x, \hat { F } _ { n } ) = \operatorname* { m i n } \biggl ( \frac { 1 } { n } \sum _ { i = 1 } ^ { n } I ( X _ { i } \leq x ), \frac { 1 } { n } \sum _ { i = 1 } ^ { n } I ( X _ { i } \geq x ) \biggr ).
\end{equation}


Consider Ton = (Ton,1,..., Ton,p) and Son = (Son,1,...,Son,p), a pair of
initial location and dispersion estimators. Here we choose for Ton and Son
respectively the coordinate-wise median and the median absolute deviation
(MAD). For each variable (Xi,, X2j,..., Xnj) (j = 1,...,P), we denote the
standardized version of Xiy by Zij = Son,j
distribution for Zj; here we use the standard normal distribution, i.e., F =
 . Let Fn, be the empirical distribution for the standardized values, that is


\begin{equation}
\hat { F } _ { n, j } ( t ) = \frac { 1 } { n } \sum _ { i = 1 } ^ { n } I ( Z _ { i j } \leq t ) \qquad j = 1, \dots, p.
\end{equation}


We define the proportion of fagged outliers by


dn,j = max(sup \{dHs(t,Fn,i)-dHs(t,F)\}+; sup \{dHs(t,Fn,j)-dHs(t,F )\}+
t -n3,j t n3.i


where n3,j = F-1( ) is a large quantile of F,. Note that, according to (),
) > () :   \} = ()g s    
which results in the simpler form written above considering the definition of
the half-space depth in the univariate case. Here, if we consider the order
statistics Z(i),j, define i- = min\{i : Z(),j > -np,j\} and i+ = max\{i : Z(i),j <
n3,\}. Using the definition of half-space depth function in the univariate case,
presented above, the previous expression can be written as


\begin{equation}
d _ { n, j } = \operatorname* { m a x } \biggl ( \operatorname* { s u p } _ { i < i _ { - } } \{ \frac { i } { n } - F _ { j } ( Z _ { ( i ), j } ) \} ^ { + }, \operatorname* { s u p } _ { i > i _ { + } } \{ F _ { j } ( Z _ { ( i ), j } ) - \frac { i - 1 } { n } \} ^ { + } \biggr ). \qquad ( 1 )
\end{equation}


Then, we fag [ndn,   observations with the smallest depth value as cell-wise
outliers and replacethembyNA's.


\subsection*{0.1  A consistent univariate, bivariate and p-variate fil
ter}


Given a sample X1,...,Xn where X; E RP,i = 1, ...,n, we first apply the
univariate filter described in the previous example to eachvariable separately.
Filtered data are indicated through an auxiliary matrix U of zeros and ones,

% [footer: 1]
\end{document}
